\subsection{Dense Video Caption}
While our dataset includes text descriptions from original webpages for many images and videos, they are often too simplistic to convey detailed visual content.
DALL-E 3~\citep{BetkerImprovingIG} has demonstrated that the ability of visual generation models to follow prompts can be significantly enhanced by training on highly descriptive generated visual captions.
Therefore, we develop an internal caption model to generate dense captions for each image and video of our dataset.
To train this model, we incorporate both open-source vision-language datasets and additional data collected in-house.

\subsubsection{Open Source Dataset}
\label{sec:open_src_cap_data}
We collect widely used vision-language datasets for both images and videos.
Our collection includes not only caption datasets but also visual Q\&A datasets focused on visual content, such as actions, counting, and OCR.
In some scenarios, we require the caption model to generate captions in a specific style or content based on user instructions.
Therefore, we also collected pure text instruction data to enhance our model's ability to follow instructions.

\subsubsection{In-house Dataset}
\label{sec:in_house_cap_data}
We collect an in-house dataset for various tasks to enhance the model's capabilities in specific areas.

\emph{Celebrities, landmarks, and movie characters.}
To train our model to recognize celebrities, landmarks, and characters, we collect dataset comprising thousands of identities.
First, we use large language models (LLMs) to collect names, then retrieve corresponding images from our image-text database using a CLIP-style model.
We test various CLIP-style models and find that TEAM~\citep{TEAM2022MM} excels in recognizing individuals, particularly Chinese celebrities.
To further reduce noise in our dataset, we implemented keyword matching for the descriptions of the image-text pairs.
An image is retained only if it is retrieved by the TEAM model and matches the specified keywords.

\emph{Object counting.}
To enhance our model's visual counting ability, we compiled a counting dataset by retrieving images with text descriptions that include words like ``one," ``two," ``three", etc.
These numbers provide coarse annotations for the images.
We then use a LLM to extract (\emph{category}, \emph{count}) pairs from the corresponding text.
Finally, we use Grounding DINO~\citep{groundingdino} to count objects in those categories.
An image is retained only if its text description matches the results from Grounding DINO~\citep{groundingdino}.

\emph{OCR.}
To improve our model's capability in generating text descriptions, we collect an OCR-augmented image caption dataset.
Initially, we used an off-the-shelf OCR detector to extract text from images.
Then, we prompt our caption model to generate descriptions based on these OCR results.
This approach ensures accurate text descriptions by using the extracted OCR text as prior knowledge.
Currently, we include only English and Chinese texts.

\emph{Camera angle and motion.}
We observe that current multi-modal language models (MLLMs) struggle with predicting camera angles and motion, including state-of-the-art commercial models like GPT-4o and Google Gemini Pro.
To enhance these predictions, we first annotate a set of videos for camera angle and motion.
This dataset can be utilized in two ways:
directly training our caption model used in the final training stage,
or by training an expert model to annotate more video data, aiding our model during the early training phase.
This dataset enhances the accuracy of our motion and angle descriptions, thereby improving the camera motion controllability of our video generation model.

\emph{Fine-grained categories.}
To enhance our model's ability to recognize fine-grained categories such as animals, plants, and vehicles, we created a dataset containing several million images across these categories.

\emph{Relational understanding.}
We enhance our model’s relational understanding by gathering dataset that focus on spatial relationships, such as left, right, up, and down, through the incorporation of existing object detection datasets.

\emph{Re-caption.}
An important capability of a caption model is to leverage existing tags or brief captions to generate dense captions.
To facilitate this, we curate a re-caption dataset that extends brief text tags or captions to detailed descriptions based on image content.

\emph{Editing instruction caption.}
We collect a dataset that describes the differences or transformations between two images.
This type of caption is valuable for image editing tasks.

\emph{Group image description.}
We collect a dataset that provides descriptions for groups of images.
Each caption begins with a description of the common features shared by the images, followed by individual descriptions of each image, using similar phrasing.

\emph{Human-annotated image and video captions.}
We gather human-annotated dense caption for both images and videos.
These represent our highest-quality caption data and are used in the final stage of model training.

\begin{figure}[htbp]
    \centering  
    \includegraphics[width=0.9\textwidth]{./figures/bench/cap_eval.pdf}  
    \caption{Our caption model achieves overall performance comparable to Google Gemini 1.5 Pro. In particular, it excels in areas such as events, camera angle, camera motion, style, and color.} 
    \label{fig:cap_eval}
\end{figure}


\subsubsection{Model Design}
\textbf{Architecture.} Our caption model adopts a LLaVA-style architecture~\citep{llava}.
We use a ViT encoder to extract visual embeddings of images and video frames.
These embeddings are projected through a two-layer perception and then fed into the Qwen LLM~\citep{qwen2.5}.

For the image input, we adopt dynamic high resolution as in LLaVA.
An image can be divided into a maximum of seven patches.
For each patch, we adaptively pool the vision embeddings to a 12$\times$12 grid representation to reduce computation.

For the video, we sample 3 frames per second, with an upper limit of 129 frames.
To further reduce computation, we employ a slow-fast encoding strategy: maintaining the original resolution for every fourth frame and applying global average pooling to the embeddings of the remaining frames.
Experiments on long video benchmarks indicate that the slow-fast mechanism enhances understanding with a limited number of visual tokens (\textit{e.g.}, improving performance from 67.6\% to 69.1\% on VideoMME~\citep{videomme}) without subtitles.

\textbf{Training.} Following recent state-of-the-art methods, our training process is divided into three stages.
In the first stage, we freeze the ViT and LLM, and train only the multi-layer perceptron to align the visual embeddings with the LLM input space, using a learning rate of 1e-3.
In the second stage, all parameters are made trainable.
In the final stage, we conduct end-to-end training on a small set of high-quality data.
For the last two stages, the learning rates are set to 1e-5 for the LLM and MLP, and 1e-6 for the ViT.

\subsubsection{Evaluation}
Automatic evaluation is crucial for the development of our caption model, allowing us to identify the strengths and weaknesses of each model version.
Following CAPability~\citep{liu2025good}, we develop an automatic caption evaluation pipeline.
We focus on ten key visual dimensions in video generation: action, camera angle, camera motion, object category, object color, object counting, OCR, scene, style, and event.
For each dimension, we randomly sampled 1,000 videos and their corresponding captions generated by our model.
We also used Google Gemini 1.5 Pro to generate dense video captions for this dataset.
We evaluate our captions and Gemini-generated captions by calculating the F1 metrics for each dimension following CAPability~\citep{liu2025good}.
The results are illustrated in Fig.~\ref{fig:cap_eval}.
Our caption performs better in video event, camera angle, camera motion, style, and object color, while Gemini excels in action, OCR, scene, object category, and object counting.